\subsection{MetalGeometryPrior: first-class constraint}
\label{sec:metal-geometry-prior}

Section~\ref{sec:mlc-types} introduced the $\Phi$-refinement type as
the type-system predicate that governs every metal combinator
$\mathsf{C}_{m}^{\Pi_{m},\,\Phi_{m}}$. This subsection unpacks the
predicate and the
\texttt{MetalGeometryPrior} runtime component that enforces it: it
fixes the geometry enumeration, lists the per-metal coordination
numbers, distinguishes dative from covalent bonds, and describes
the MCTS scoring integration that prunes geometry-violating terms
before $\beta$-reduction.

\paragraph{Honest framing.}
The geometry enumeration and per-metal coordination numbers are
\textsc{measured}: they are imported verbatim from
\texttt{molmetal/molmetal\_lam/priors/metal\_geometry.py}. The
$\Phi$-refinement is \textsc{measured} on the round-10 per-component
harness across four canonical metals
(\texttt{round10\_per\_component\_metrics.md}, row 2).
The MCTS pruning integration is \textsc{measured} on the
WF-Lambda-1c pilot (\texttt{wf\_lambda1c\_pilot\_v3/final.md}).
Broader multi-metal coverage (Fe, Mn, Cu, Pd, Rh, Os) is
\textsc{projected} future work (Section~\ref{sec:future}).

\subsubsection{Coordination geometry as a type-system predicate}
\label{sec:metal-geom-enum}

The metal refinement type
(Eq.~\eqref{eq:metal-refinement}) requires a metal combinator to
satisfy two projections: an integer coordination number $n_m$ and
a literal geometry tag $g_m$ drawn from the four-element
enumeration
\begin{equation}
  \label{eq:geometry-enum}
  \mathrm{Geometry} \;=\; \{
    \textsc{SquarePlanar},\;
    \textsc{Octahedral},\;
    \textsc{Tetrahedral},\;
    \textsc{TrigonalBipyramidal}
  \}.
\end{equation}
Equivalently, for a metal atom with atomic number $Z$ in the
shipped registry, the typed combinator must satisfy
\begin{equation}
  \label{eq:coordination-geom-predicate}
  \mathrm{coordination\_geometry}(\mathsf{C}_{m}^{\Pi_{m},\,\Phi_{m}})
    \;:\; \mathcal{V} \times \mathbb{M}
      \;\longrightarrow\; \mathrm{Geometry},
\end{equation}
and the refinement $\Phi_{m}$ rejects any term whose realised
geometry does not match the lookup.

\paragraph{Source.}
The predicate is materialised in
\texttt{molmetal/molmetal\_lam/priors/metal\_geometry.py}, where
\texttt{DEFAULT\_METAL\_GEOMETRY}
(\texttt{metal\_geometry.py:364--374}) is a hand-curated
atomic-number-to-geometry table:
\begin{verbatim}
DEFAULT_METAL_GEOMETRY: dict = {
    78: Geometry.SQUARE_PLANAR,      # Pt(II)
    46: Geometry.SQUARE_PLANAR,      # Pd(II)  (analogous d8)
    29: Geometry.TETRAHEDRAL,        # Cu(I) / Zn(II) common coordination
    30: Geometry.TETRAHEDRAL,        # Zn(II)
    26: Geometry.OCTAHEDRAL,         # Fe(II)
    44: Geometry.OCTAHEDRAL,         # Ru(II)
    77: Geometry.OCTAHEDRAL,         # Ir(III)
    45: Geometry.OCTAHEDRAL,         # Rh(III)
    25: Geometry.TRIGONAL_BIPYRAMIDAL,  # Mn(0)
}
\end{verbatim}
The literal type \texttt{Geometry(str)} is a closed string-enum
that the prior dispatches on without importing the \texttt{enum}
stdlib module (\texttt{metal\_geometry.py:341--350}); the
\texttt{SUPPORTED\_GEOMETRIES} tuple
(\texttt{metal\_geometry.py:355--360}) is the constant set used by
the asserters in \texttt{prior\_loss}.

\subsubsection{Per-metal coordination numbers}
\label{sec:metal-coord-numbers}

The five metal combinators currently registered in
\texttt{METAL\_ATOMS} (\texttt{atoms/combinators.py:195--225})
carry the refinement pairs listed in
Table~\ref{tab:metal-types-3-3}, which restate
Table~\ref{tab:metal-types} with the lookup tags needed by the
prior.

\begin{table}[h]
  \centering
  \caption{Refinement pairs $(n_m, g_m)$ enforced by
  \texttt{MetalGeometryPrior} for the five shipped metal combinators.
  $Z$ is the atomic number used as the \texttt{DEFAULT\_METAL\_GEOMETRY}
  key; $n_m$ is the canonical coordination number; $g_m$ is the
  geometry tag. \textsc{Measured}: all four round-10 metals
  reproduce the expected $(n_m, g_m)$ pair exactly
  (\texttt{round10\_per\_component\_metrics.md},
  METAL\_GEOMETRY\_OK\,=\,1.000).}
  \label{tab:metal-types-3-3}
  \begin{tabular}{lcccc}
    \hline
    Metal combinator & $Z$ & Oxidation & $n_m$ & $g_m$ \\
    \hline
    $\mathsf{C}_{Pt_{II}}^{\Pi_{Pt},\,\Phi_{Pt}}$   & 78 & $+2$ & 4 & square-planar \\
    $\mathsf{C}_{Ru_{II}}^{\Pi_{Ru},\,\Phi_{Ru}}$   & 44 & $+2$ & 6 & octahedral    \\
    $\mathsf{C}_{Zn_{II}}^{\Pi_{Zn},\,\Phi_{Zn}}$   & 30 & $+2$ & 4 & tetrahedral   \\
    $\mathsf{C}_{Ir_{III}}^{\Pi_{Ir},\,\Phi_{Ir}}$  & 77 & $+3$ & 6 & octahedral    \\
    $\mathsf{C}_{Au_{III}}^{\Pi_{Au},\,\Phi_{Au}}$  & 79 & $+3$ & 4 & square-planar \\
    \hline
  \end{tabular}
\end{table}

\paragraph{Measured (round-10 per-component harness).}
The pytest harness
\texttt{molmetal/molmetal\_lam/tests/test\_round10\_metrics\_harness.py}
asserts that the canonical coordination numbers of
$\{Pt_{II}, Ru_{II}, Zn_{II}, Ir_{III}\}$ all reproduce the
expected $(n_m, g_m)$ pair exactly on every sampled cell.
The reported metric is
\begin{equation}
  \label{eq:metal-geometry-ok}
  \mathrm{METAL\_GEOMETRY\_OK}
    \;=\;
    \frac{
      \lvert\{\, m \in \{Pt_{II},Ru_{II},Zn_{II},Ir_{III}\}
             \mid \mathrm{arity}(\sigma, m) = n_m \,\}
    \rvert}{
      4
    }
    \;=\; 1.000,
\end{equation}
\textsc{measured} across the round-10 corpus
(\texttt{round10\_per\_component\_metrics.md}, row 2: $4/4$
matches, no fallbacks to the generic
\texttt{Atom} constructor). The Round-3
Au$_{\text{III}}$ arity property-test bug does not regress.

\subsubsection{Dative bonds and the FreeSiteLedger distinction}
\label{sec:dative-bonds}

A dative bond is the first-class construct that the MLC uses to
represent metal--ligand coordination: both electrons of the bond
are supplied by a single donor atom, so the bond is \emph{directed}
and \emph{asymmetric} in a way that a covalent bond is not. In the
$\Pi$-type of Eq.~\eqref{eq:pi-type} the asymmetry is captured by
the constraint
\begin{equation}
  \label{eq:dative-asymmetry}
  b = \texttt{dative}
    \;\Longrightarrow\;
    a.\omega > 0 \;\wedge\; d.\omega \leq 0,
\end{equation}
i.e.\ the acceptor $a$ must have a non-zero formal oxidation state
(an electron-poor centre) while the donor $d$ must be
electron-rich (or neutral with a lone pair).

\paragraph{Two bond kinds, one ledger.}
The runtime distinguishes the two kinds at the
\texttt{Bond} factory level
(\texttt{bonds/application.py:135--138}):
\begin{verbatim}
COVALENT = "covalent"
DATIVE   = "dative"
AROMATIC = "aromatic"
HYDROGEN = "hydrogen"
\end{verbatim}
A \texttt{Bond} with \texttt{kind = DATIVE} carries the
\texttt{is\_dative = True} flag (\texttt{bonds/application.py:537})
and is recorded in a separate counter in the
\texttt{FreeSiteLedger} (\texttt{bonds/application.py:243--325}):
the acceptor loses one free site (currying), while the donor's
free-site count collapses to zero via
\texttt{AtomSite.saturate()}. Covalent and aromatic bonds, by
contrast, debit \emph{both} partners by the bond order. The
\texttt{FreeSiteLedger.check\_arity\_conservation} method
(\texttt{bonds/application.py:284--325}) is the runtime check
that the dative consumption rule balances; a candidate term that
violates it is rejected before $\beta$-reduction.

\paragraph{Measured (round-10 per-component harness).}
The reported metric is
\begin{equation}
  \label{eq:dative-fraction}
  \mathrm{DATIVE\_FRACTION}
    \;=\;
    \frac{
      \lvert\{\, b \in \mathrm{bonds}
             \mid b.\mathrm{kind} = \texttt{DATIVE} \,\}
    \rvert}{
      \lvert\{\, b \in \mathrm{bonds} \,\}\rvert
    }
    \;=\; 1.000,
\end{equation}
\textsc{measured} on the round-10 corpus: every bond minted via
\texttt{Bond.dative} against a metal atom lands in
\texttt{DATIVE}, never in \texttt{COVALENT}
(\texttt{round10\_per\_component\_metrics.md}, row 3; 20
coordination bonds sampled, 20 dative, 0 covalent). The
\texttt{BOND\_KIND\_DISTRIBUTION} row reports all four kinds
present in non-degenerate counts, so the distinction is not
degenerate.

\subsubsection{MCTS scoring integration}
\label{sec:mcts-geometry-scoring}

The prior is wired into the Lambda MCTS loop at
\texttt{molmetal/scripts/r4\_lambda\_only\_run.py:236--267} as a
scoring bonus that fires \emph{before} $\beta$-reduction:
\begin{verbatim}
def metal_geometry_prior_bonus(state, *, enabled, ...) -> float:
    """+1.0 if the molecule satisfies the metal-geometry prior."""
    if not enabled: return 0.0
    ...
    has_metal, matches_target = ..., ...
    return 1.0 if (has_metal and matches_target) else 0.0
\end{verbatim}
The bonus contributes $+1.0$ to the per-state MCTS score when the
candidate (a) contains at least one metal atom from the metal
alphabet $\{\,Pt, Ru, Ir\,\}$ (extensible to the full
$\mathbb{M}$ in Table~\ref{tab:metal-types-3-3}), and (b) the
metal atom's realised coordination number matches the
\texttt{DEFAULT\_METAL\_GEOMETRY} target. When \texttt{enabled} is
\texttt{False} the prior is a strict no-op (returns 0.0) --- this
is the Round-10 axis-D ``prior OFF'' ablation path.

\paragraph{Pruning semantics.}
A term that violates the $\Phi$-refinement is pruned \emph{before}
the $\beta$-reduction engine is invoked. This is what makes the
$80\%$ drop in candidate count in the
Section~\ref{sec:click-chem-ablation} ablation: the five-rule
search \emph{generates} more candidates than CuAAC alone, but the
MetalGeometryPrior rejects the geometrically incompatible ones.
The ablation in
Section~\ref{sec:click-chem-ablation} therefore measures
post-filter candidate count, not raw generation count, and the
honest-framing note in that subsection makes the distinction
explicit.

\paragraph{Pipeline view (Fig 2).}
Figure~\ref{fig:pipeline} shows the prior as stage~(c) on the
main spine of the pocket-conditioned Lambda pipeline. The dashed
arrow from stage~(d) (RDKit decode) to stage~(c) (geometry
prior) is the feedback loop that re-runs the prior on every
assembled candidate before the validity gate at stage~(d) emits
the SMILES to the parallel validation band (stages~(e)--(h)) and
the \texttt{RewardAggregator} at stage~(i).

\subsubsection{Empirical evidence from WF-Lambda-1c}
\label{sec:metal-geom-evidence}

The WF-Lambda-1c pilot
(\texttt{molmetal/reports/wf\_lambda1c\_pilot\_v3/final.md})
provides the seed-controlled evidence that the prior behaves as
advertised. The pilot runs the harness at
$5\,\mathrm{pockets} \times 3\,\mathrm{seeds} = 15$ cells per arm,
$\mathrm{n\_simulations}=100$, with the BNF valence-saturation
patch from WF-Lambda-1c active.

\paragraph{Clean seed-controlled ablation.}
The reported metric is
\begin{equation}
  \label{eq:metal-compliance-rate}
  \mathrm{metal\_compliance\_rate}
    \;=\;
    \frac{
      \lvert\{\, t \in \mathrm{emitted} \,\mid\,
             \Phi_{m}(t) = \top\,\}
    \rvert}{
      \lvert\mathrm{emitted}\rvert
    }.
\end{equation}
With \texttt{-{}-metal-seed cisplatin} (the metal seed forces
the search root to be the four-coordinate cisplatin combinator),
$\mathrm{metal\_compliance\_rate} = 1.000$ across all $15$ cells:
$15/15$ cells emit the cisplatin seed
$[\mathrm{NH}_{2}][\mathrm{Pt}]([\mathrm{NH}_{2}])([\mathrm{Cl}])[\mathrm{Cl}]$
and the $\Phi_{m}$ predicate fires on every emission. Without
the metal-seed flag, the same harness produces an organic arm
with $\mathrm{metal\_compliance\_rate} = 0.000$: the roots are
organic and no Pt/Ru/Ir centre appears. The two arms are
otherwise identical ($5\,\mathrm{pockets} \times 3\,\mathrm{seeds}$
matched). The $\Delta = 1.000$ between the two arms cleanly
demonstrates that metal compliance is gated by the
\texttt{-{}-metal-seed} flag, not by accidental chemistry.

\paragraph{All six metrics non-zero in the main arm.}
The main arm (with cisplatin seed) reports
$\mathrm{validity\_rate}=1.000$,
$\mathrm{uniqueness\_rate}=1.000$,
$\mathrm{synthesizability\_rate}=1.000$,
$\mathrm{metal\_compliance\_rate}=1.000$,
$\mathrm{novelty}=1.000$, and
$\mathrm{diversity\_alpha}=0.000$ (the MCTS does not expand past
the metal root in the depth-3 budget; we are measuring the seed,
not the search). The non-metal arm reports the same five
metrics at 1.000 except $\mathrm{metal\_compliance\_rate}=0.000$
and $\mathrm{diversity\_alpha}=0.0049$ (the organic arm sees the
canonical organic diversity). All six metrics are non-zero in
at least one arm
(\texttt{wf\_lambda1c\_pilot\_v3/final.md}, \textsc{Measured}).

\subsubsection{Honest framing: scope and limits}
\label{sec:metal-geom-honest}

The $\Phi$-refinement and \texttt{MetalGeometryPrior} have been
exercised on three seed families only:
\begin{itemize}
  \item \emph{Cisplatin family} (\texttt{-{}-metal-seed cisplatin}):
    Pt$_{\text{II}}$ square-planar, four donors
    ($2 \times \mathrm{NH}_{3}$, $2 \times \mathrm{Cl}$).
    \textsc{Measured} on $5 \times 3 = 15$ cells with full
    control ablation.
  \item \emph{Ruthenium-arene family} (\texttt{-{}-metal-seed
    ru\_arene}): Ru$_{\text{II}}$ octahedral, six donors
    ($\eta^{6}$-arene + 3 monodentate ligands). \textsc{Measured}
    in the spot-check test suite
    (\texttt{tests/test\_metal\_geometry.py}); not exercised at
    the cell level of the WF-Lambda-1c pilot.
  \item \emph{Iridium-Cp$^{\star}$ family} (\texttt{-{}-metal-seed
    ir\_cpstar}): Ir$_{\text{III}}$ octahedral, six donors
    ($\eta^{5}$-Cp$^{\star}$ + 2 monodentate ligands). \textsc{Measured}
    in the spot-check test suite; not exercised at the cell level
    of the WF-Lambda-1c pilot.
\end{itemize}
Broader metal validation --- Fe$_{\text{II}}$ octahedral,
Mn$_{\text{0}}$ trigonal-bipyramidal, Cu$_{\text{I}}$/Zn$_{\text{II}}$
tetrahedral, Pd$_{\text{II}}$ square-planar, Rh$_{\text{III}}$
octahedral, Os$_{\text{II}}$ octahedral --- is deferred to the
multi-metal extension in Section~\ref{sec:future}.
The 1c patch lifts the synthesis oracle for organic AND metal roots
but does \emph{not} constitute a guarantee at scale; a larger
sweep may surface cells where transition-metal valences are set
loosely, which is a separate audit item.

\paragraph{Why a 5 $\times$ 3 pilot is still evidence.}
The WF-Lambda-1c sample is the smallest scientifically meaningful
pilot ($n = 30$ cells, matched across two arms) and is \emph{not}
a sweep. Its purpose is to (i) verify that the geometry prior fires
deterministically on a metal-seeded root, and (ii) verify that
the prior does not spuriously fire on an organic root. Both
properties are demonstrated by the $\Delta = 1.000$ ablation.
The full 100-pocket $\times$ 3-seed sweep is
\textsc{projected} to preserve the same per-cell behaviour at the
aggregate level, but is left to future work
(Section~\ref{sec:future}, multi-metal extension).

\paragraph{Relationship to the closure theorem.}
The $\Phi$-refinement is a \emph{local} constraint on a single
metal combinator; it is independent of the
\emph{global} claim that every well-typed closed term admits a
$\beta$-NF (Section~\ref{sec:mlc-strong-norm}, Appendix~\ref{appendix:closure-theorem}).
The two interact: the closure theorem guarantees that the
$\beta$-NF operator terminates, while the
MetalGeometryPrior guarantees that the metal combinator satisfies
the refinement along the way. Together they imply that every
well-typed term whose metal subterm satisfies $\Phi_{m}$ admits a
$\beta$-NF with a closed coordination sphere. The full formal
statement is given in Appendix~\ref{appendix:closure-theorem}.

% --- cross-refs ----------------------------------------------------------
% Figure~\ref{fig:pipeline} (Section 4) cites stage~(c) on the main
% spine as the runtime hook for this prior.
% Section~\ref{sec:eval-metrics} cites METAL_GEOMETRY_OK = 1.000 as
% the per-component sanity floor.
% Section~\ref{sec:future} cites the multi-metal extension as the
% deferred scope expansion for this prior.
